\documentclass[11pt]{article}

\usepackage[T1]{fontenc}
\usepackage[utf8]{inputenc}
\usepackage{amsmath,amssymb}
\usepackage{booktabs}
\usepackage{graphicx}
\usepackage{microtype}
\usepackage{geometry}
\usepackage{xcolor}
\usepackage{hyperref}
\usepackage{enumitem}

\geometry{margin=1in}
\hypersetup{
  colorlinks=true,
  linkcolor=blue!55!black,
  citecolor=blue!55!black,
  urlcolor=blue!55!black
}
\setlist{nosep}

% Provisional title and author order: verify before submission.
\title{HayFlow: Towards a Causal, Morphology-Aware Neural Surrogate\\
for a Biophysical Cortical Neuron}
\author{Alessio Ragno \and Alessandro Belli}
\date{August 2026}

\begin{document}

\maketitle

\begin{abstract}
Detailed multicompartment neuron models reproduce nonlinear dendritic and
somatic dynamics, but their computational cost limits their use in large
network simulations and iterative experiments. Existing neural surrogates can
predict membrane voltage accurately, yet voltage-only input--output mappings do
not expose the internal variables needed to define a causal, restartable, and
autoregressive simulator. We introduce HayFlow, a developing surrogate of the
Hay layer-5 pyramidal-cell model that learns a one-millisecond flow map over the
complete 642-segment morphology, a rich boundary state, and realized synaptic
input. Our methodology decomposes the problem into controlled learnability,
representation, optimization, and recursive-exposure experiments before
scaling the model. A frozen morphology-aware core followed by a direct-tree
residual decoder achieves an ensemble RMSE of \textbf{0.3308 mV} on a
preregistered fresh one-step test, compared with 2.5763 mV for persistence.
In a separate information-matched comparison, the compact HayFlow voltage path
outperforms an 8,002-parameter Branch-ELM core in every paired seed, reducing
median global RMSE by 6.44\% and somatic RMSE by 34.94\%. However, the strongest
one-step model fails under closed-loop rollout. A later recursive prototype
shows that four-millisecond pushforward exposure reduces the median eight-step
RMSE from 18.162 mV for a support-matched passive reference to 14.720 mV, with
no physical-voltage violations. This prototype does not pass its cross-seed
promotion gate and still receives teacher mechanism state during rollout.
Thus, the current evidence establishes an accurate one-step flow map and a
causally supported route toward recursive stability, but not yet an autonomous
replacement for NEURON.
\end{abstract}

\section{Introduction}

Biophysically detailed neuron models represent membrane voltage, ion-channel
gates, calcium dynamics, synaptic conductances, axial coupling, and nonlinear
dendritic events over a branched morphology. The Hay layer-5 pyramidal-cell
(L5PC) model is a prominent example: it captures a broad range of dendritic and
perisomatic active properties and is routinely simulated with NEURON
\cite{hay2011,hines1997}. This level of mechanistic detail is valuable, but it
is expensive when the neuron must be evaluated repeatedly or embedded in a
large network.

Neural surrogates offer a promising alternative. Beniaguev, Segev, and London
showed that the input--output transformation of a detailed cortical neuron can
be approximated by a deep artificial neural network \cite{beniaguev2021}.
Nevertheless, predicting voltage from a finite history and constructing a
restartable dynamical system are different objectives. In closed loop, a
surrogate consumes its own predicted state. Small errors in voltage, slow
channel variables, calcium, or synaptic tails may then change voltage-dependent
conductances and threshold events, causing the model to enter states absent
from teacher-forced training.

HayFlow addresses this distinction explicitly. We formulate the teacher as a
causal one-millisecond flow map over a rich state contract, preserve the
instantiated morphology, and separate three questions:

\begin{enumerate}
  \item Is the one-step transition identifiable from the available causal
  state and input?
  \item Can a compact, morphology-aware model learn that transition?
  \item Does the learned transition remain stable when recursively composed?
\end{enumerate}

The present work is a technical progress report rather than a claim of a
finished simulator. Its contributions are:

\begin{itemize}
  \item a replay-validated teacher and dataset contract covering all 642
  instantiated segments, including soma, axon initial segment, axon, and
  dendrites;
  \item a decision-grade one-step checkpoint family evaluated on a sealed,
  preregistered fresh test;
  \item an information-matched comparison showing a modest but consistent
  advantage over a compact Branch-ELM baseline under equal numerical input;
  \item a controlled diagnosis showing that excellent one-step accuracy does
  not imply autoregressive stability;
  \item factorial evidence that pushforward boundary exposure, rather than the
  tested ordered-event auxiliary or directional penalty, materially improves
  learned recursive dynamics.
\end{itemize}

\section{Related Work}

\paragraph{Biophysical neuron models.}
The Hay L5PC model combines detailed morphology with active conductances tuned
to reproduce dendritic and perisomatic behavior \cite{hay2011}. NEURON provides
the numerical environment used to instantiate and integrate such
multicompartment models \cite{hines1997}. HayFlow does not modify this teacher:
the original NEURON model remains the canonical source of trajectories,
internal states, and event labels.

\paragraph{Neural surrogates of cortical neurons.}
The NeuronIO work demonstrated that a detailed cortical neuron can be
approximated by a deep temporal convolutional network and also introduced a
small branch-structured baseline \cite{beniaguev2021}. That work primarily
models the input--output transformation from presynaptic events to somatic
voltage and spikes. HayFlow instead targets a causal boundary-to-boundary
transition over the complete instantiated state and morphology. Consequently,
published error values cannot be compared directly unless support, target,
information, and metric are aligned.

\paragraph{Learned physical simulation.}
Graph-based learned simulators exploit relational structure to model physical
systems and meshes \cite{sanchez2020,pfaff2021}. HayFlow shares the view that
topology should be represented explicitly, but the neuronal problem adds stiff
local mechanism states, receptor dynamics, voltage-dependent nonlinearities,
and threshold events. We therefore retain a passive Hines-style physical prior
in the recursive line of experiments rather than asking a generic graph model
to rediscover all cable coupling from data.

\paragraph{Autoregressive distribution shift.}
Teacher forcing creates a mismatch between training on ground-truth states and
deployment on model-generated states. Scheduled sampling and learned-simulator
noise curricula address related exposure problems \cite{bengio2015,
sanchez2020}. Our pushforward intervention follows the same general principle,
but is evaluated causally within a paired factorial design and against a
support-matched passive reference.

\section{Problem Formulation and Teacher Contract}

Let $S_t$ denote the complete causal teacher state at a millisecond boundary
and let $U_{[t,t+1)}$ denote the realized inputs during the following interval.
The target flow map is
\begin{equation}
  \mathcal{F}_{1\,\mathrm{ms}}:
  \left(S_t,U_{[t,t+1)}\right)
  \longmapsto
  \left(S_{t+1},V_{[t,t+1]},E_{[t,t+1)}\right),
  \label{eq:flowmap}
\end{equation}
where $V_{[t,t+1]}$ is a high-resolution voltage microtrace on selected probes
and $E_{[t,t+1)}$ contains event labels such as somatic spikes,
backpropagating action potentials, calcium spikes, and NMDA events.

The state contract contains, where available in the authentic teacher:

\begin{itemize}
  \item voltage for every one of the 642 instantiated segments;
  \item ion-channel and other NMODL \texttt{STATE} variables;
  \item calcium and accessible ionic concentrations;
  \item synaptic rise/decay state and causal release information;
  \item morphology, mechanism-presence masks, and stable segment indices;
  \item random-stream metadata needed for reproducibility.
\end{itemize}

The realized input $U_{[t,t+1)}$ records the ordered synaptic events, their
intra-millisecond timestamps, synapse identities, receptor-aware effects, and
somatic current when present. Realized release is causal in this formulation:
the authentic synaptic front-end determines release before the membrane update,
so it is an admissible input to the dynamical core rather than information from
the future endpoint.

Teacher snapshots are restored and replayed under identical inputs. The
diagnostic dataset passed exhaustive replay with a maximum discrepancy below
$10^{-5}$ in its registered validation. This contract makes it possible to
distinguish model error from teacher nondeterminism or indexing drift.

\section{Method}

\subsection{Component-Level Experimental Methodology}

Instead of immediately scaling a large end-to-end architecture, we use a
component-level experimental workflow. Each experiment isolates one proposed
capability---for example state update, voltage coupling, axial transport,
event use, or recursive exposure---and includes an aligned negative control.
When several hypotheses share data, seeds, and minibatches, they are evaluated
in a small preregistered factorial matrix. Checkpoints at increasing budgets
provide miniature scaling curves without duplicating training.

The workflow distinguishes four failure modes:

\begin{enumerate}
  \item \textbf{information}: the causal input does not identify the target;
  \item \textbf{representation}: the useful information is present but poorly
  encoded;
  \item \textbf{optimization}: the representation is expressive but not used
  by the optimizer;
  \item \textbf{recursive exposure}: a one-step learner is evaluated on states
  generated by its own imperfect transition.
\end{enumerate}

Probe tasks, target permutations, path deletions, paired seeds, and frozen
historical checkpoints are used to discriminate these explanations. Test and
fresh-test outcomes are not used for architecture or checkpoint selection.

\subsection{Model A: One-Step HayFlow-Hines Decoder}

The strongest one-step family consists of a frozen HayFlow-Hines H2 core and a
trainable direct-tree residual decoder. The decoder receives 226 fixed
multiscale tree features per segment, concatenated with a learned
16-dimensional segment embedding. A shared SiLU multilayer perceptron of width
$242\rightarrow96\rightarrow96\rightarrow1$ predicts a bounded voltage
residual:
\begin{equation}
  \widehat{V}_{t+1,i}=V^{\mathrm{H2}}_{t+1,i}
  +120\tanh\!\left(g_\theta(\phi_{t,i},e_i)/120\right),
\end{equation}
where $i$ indexes a segment, $\phi_{t,i}$ denotes its tree features, and $e_i$
is its learned embedding. The decoder has 43,009 trainable parameters; the
frozen H2 parameters are additional and are not included in this count.

Three fixed seeds were fitted on 120 train-only boundary pairs. Thirty
disjoint internal pairs were used for checkpoint selection, and previously
opened development support was evaluated only after freezing. A separate
32-pair, 64-episode fresh-test plan was then generated and opened without
retraining or seed selection.

\subsection{Model B: Recursive Physical-Prior Prototype}

The current learned recursive frontier uses a passive Hines cable solve as a
physical prior and a 12,972-parameter neural cell that predicts an effective
membrane source from causal state and compact input. Its strongest arm uses
four-millisecond pushforward exposure: during training, the model is exposed
to boundary states generated by its own preceding predictions before the loss
is evaluated. This directly targets the deployment distribution.

The registered $2\times2\times2$ design crossed:

\begin{itemize}
  \item a causal ordered-event auxiliary objective on/off;
  \item teacher-boundary versus four-millisecond pushforward exposure;
  \item a passive-relative directional stability penalty on/off.
\end{itemize}

All arms shared initialization, windows, minibatches, and budget. Calibration
selection and arm selection used disjoint calibration halves. The development
role was opened only after selection.

\section{Experimental Setup}

The base diagnostic resource contains 29,240 transitions and is augmented by a
640-transition validation top-up, giving 29,880 transitions without physically
merging the stores. The data are split by seed and protocol rather than by
random windows from the same trajectory. A separate 05j-o fresh teacher shard
contains 768 transitions from 64 episodes arranged as 32 paired examples.

The principal voltage metric is
\begin{equation}
  \mathrm{RMSE}_V =
  \sqrt{\frac{1}{N T}\sum_{n=1}^{N}\sum_{i=1}^{T}
  \left(\widehat{V}_{n,i}-V_{n,i}\right)^2},
\end{equation}
reported in millivolts. We additionally monitor region-specific and
activity-conditioned RMSE, maximum segment error, mean drift, branching
retention, and physical-voltage violations. One-step and recursive metrics are
reported separately and must not be merged into a single leaderboard.

\subsection{Information-Matched Branch-ELM Comparison}

To answer whether the apparent advantage over the small Branch-ELM baseline
was caused by the enriched data contract rather than the architecture, we ran
a corrective information-matched comparison. Both models receive the exact
same 76-value per-segment numerical tensor, containing boundary voltage and
axial differences, mechanism state and presence, local ions, realized causal
input, static segment features, and region identity. They use the same
transitions, sample order, loss, optimizer, and unclipped target
$V_{t+1}-V_t$. Neither model receives $V_{t+1}$ as input, and no rollout or
spike metric is part of this comparison.

The Branch-ELM core has 8,002 trainable parameters (7,981 on its voltage
influence path). The HayFlow voltage bridge has 8,985 trainable parameters.
Its downstream 7,212-parameter state updater is frozen and cannot affect the
scored voltage. The active paths are therefore close, although not exactly
parameter matched. Training uses three paired seeds, 1,024 fit transitions,
256 calibration transitions, and 256 disjoint development transitions per
registered execution.

\section{Results}

\subsection{Fresh One-Step Evaluation}

Table~\ref{tab:fresh} reports the preregistered 05j-o fresh test. All three
frozen seeds pass the registered gate and improve substantially over both
persistence and frozen H2. The canonical report result is the three-seed
ensemble, not a seed chosen after opening the fresh test.

\begin{table}[t]
  \centering
  \caption{Preregistered one-step fresh-test voltage RMSE. The ensemble is the
  canonical Model A result.}
  \label{tab:fresh}
  \begin{tabular}{lrr}
    \toprule
    Model & RMSE (mV) & Status \\
    \midrule
    Persistence & 2.5763 & Baseline \\
    Frozen HayFlow-Hines H2 & 2.7445 & Baseline \\
    Model A, seed 17 & 0.4037 & Passed \\
    Model A, seed 29 & 0.3774 & Passed \\
    Model A, seed 43 & 0.3775 & Passed \\
    Model A, three-seed ensemble & \textbf{0.3308} & Authorized one-step \\
    \bottomrule
  \end{tabular}
\end{table}

The ensemble reduces RMSE by 87.16\% relative to persistence. Its authorization
is strictly one-step: it does not establish stable recursive simulation.

\subsection{Information-Matched Comparison with Branch-ELM}

Table~\ref{tab:elm} confirms the qualitative observation discussed during the
project review: under equal numerical information, HayFlow is slightly better
globally and more clearly better in somatic and active regimes. HayFlow wins
the global paired comparison for every seed; the per-seed global reductions
are 3.15\%, 6.44\%, and 9.96\%.

\begin{table}[t]
  \centering
  \caption{Information-matched one-step development comparison. Values are
  paired-seed median RMSE. This is not a rollout or fresh-test comparison.}
  \label{tab:elm}
  \begin{tabular}{lrrr}
    \toprule
    Scope & Branch-ELM & HayFlow & Reduction \\
          & RMSE (mV) & RMSE (mV) & (\%) \\
    \midrule
    All 642 segments & 6.380 & \textbf{5.995} & 6.44 \\
    Soma & 14.824 & \textbf{9.577} & 34.94 \\
    Active examples & 21.012 & \textbf{18.867} & 8.92 \\
    \bottomrule
  \end{tabular}
\end{table}

This comparison supports the compact causal and morphology-aware HayFlow
voltage representation. It does not justify directly comparing the 0.3308 mV
fresh-test ensemble with the published Branch-ELM value: those figures use
different target scopes, supports, and metrics. The matched comparison is the
scientifically valid answer to the architecture question.

\subsection{One-Step Accuracy Does Not Compose Automatically}

When the frozen Model A checkpoints are repeatedly applied to their own
predicted state, their eight-millisecond RMSEs rise to 76.70, 95.47, and
62.72 mV. New synaptic inputs at each millisecond do not reset this error:
voltage-dependent currents, slow mechanism state, calcium, and synaptic tails
carry previous state errors forward. The result rules out treating one-step
RMSE as a sufficient proxy for deployment stability.

\subsection{Recursive Progress from Pushforward Exposure}

The 06b-r factorial experiment identifies pushforward exposure as the only
materially supported intervention among the three tested factors. The best
learned arm reaches a median 8 ms development RMSE of 14.720 mV, compared with
18.162 mV for the support-matched passive reference and 46.938 mV for the
previous frozen learned arm on the same support. This corresponds to an
18.95\% reduction relative to the passive reference and a 68.64\% reduction
relative to the previous learned arm. It produces zero physical-voltage
violations.

\begin{table}[t]
  \centering
  \caption{Registered causal readout of the 06b-r factorial experiment.}
  \label{tab:factorial}
  \begin{tabular}{lrl}
    \toprule
    Intervention or control & Marginal effect & Interpretation \\
    \midrule
    Pushforward exposure, 4 ms & 18.96\% gain & Supported \\
    Ordered-event auxiliary & 0.0058\% gain & Inert at tested scale \\
    Directional stability penalty & 0.00004\% gain & Inert at tested scale \\
    Delete ordered-event path & 0.0145\% error change & Not material \\
    Delete all compact causal $U$ & 2.17\% error increase & Material input path \\
    \bottomrule
  \end{tabular}
\end{table}

The arm is not promoted. Seed 61043 regresses by 10.08\% on calibration half B,
violating the preregistered no-seed-regression gate. In addition, the model is
supplied with teacher mechanism state during rollout and is therefore not yet
an autonomous surrogate. Development results cannot be used retrospectively
to overturn the calibration veto.

\section{Discussion}

The experiments provide two complementary results. First, the causal boundary
state makes the one-millisecond transition highly predictable: Model A reaches
sub-millivolt fresh-test error and clearly beats simple one-step baselines.
Second, recursive composition is a distinct problem. The one-step champion can
leave the teacher-state manifold after only a few predictions, at which point
small errors are amplified by nonlinear membrane mechanisms.

The information-matched Branch-ELM experiment narrows the architectural claim.
When both compact models receive identical causal information, HayFlow has a
consistent but moderate global advantage and a larger somatic advantage. This
suggests that morphology-aware voltage features are useful, but it does not
show that model identity alone explains the one-step champion: Model A is
larger, includes a frozen H2 core, and is evaluated under a separate fresh-test
contract.

The recursive factorial result is encouraging because it separates a real
repair from several plausible but ineffective additions. The tested event
auxiliary and directional penalty did not improve the selected metric, while
pushforward exposure did. Thus, the immediate bottleneck is not simply a lack
of nominal network capacity or a missing event feature. Training-distribution
alignment is causal and material. At the same time, cross-seed fragility and
teacher-state dependence show that exposure repair alone is insufficient.

\subsection{Limitations}

\begin{itemize}
  \item The current datasets are intentionally diagnostic and enriched; they
  do not yet cover the complete biological input distribution at the scale of
  the original NeuronIO study.
  \item The canonical 0.3308 mV result is one-step only.
  \item The recursive 14.720 mV result is development evidence from a model
  that failed its calibration promotion gate.
  \item Model B receives teacher mechanism state during rollout, preventing an
  autonomous deployment claim.
  \item The Branch-ELM comparison is train-derived development evidence, not a
  new sealed test, and the active voltage paths differ modestly in parameter
  count.
  \item Spike timing, rapid firing, long plateaus, longer horizons, and network
  integration remain to be validated.
\end{itemize}

\section{Conclusion and Future Work}

HayFlow currently combines a fresh-tested 0.3308 mV one-step checkpoint family
with a separate pushforward-trained recursive prototype that improves 8 ms
error by 18.95\% over its support-matched physical baseline. The matched
Branch-ELM comparison further shows that, with identical causal inputs, the
compact HayFlow voltage representation is consistently better, although only
modestly so in the global metric. These results constitute meaningful progress
toward a morphology-aware neural simulator, but they do not yet establish an
autonomous replacement for NEURON.

The next experiments should retain the passive physical prior, compact causal
input, and pushforward exposure while removing causally inert branches. They
should then test nested cross-seed calibration, autonomous mechanism-state
prediction, multiple pushforward horizons and budgets, event-positive regimes,
and longer closed-loop rollouts. A new sealed test should be opened only after
the resulting candidate passes all train-derived and calibration gates.

\begin{thebibliography}{9}

\bibitem{hay2011}
E. Hay, S. Hill, F. Sch\"urmann, H. Markram, and I. Segev.
Models of neocortical layer 5b pyramidal cells capturing a wide range of
dendritic and perisomatic active properties.
\emph{PLoS Computational Biology}, 7(7):e1002107, 2011.

\bibitem{hines1997}
M. L. Hines and N. T. Carnevale.
The NEURON simulation environment.
\emph{Neural Computation}, 9(6):1179--1209, 1997.

\bibitem{beniaguev2021}
D. Beniaguev, I. Segev, and M. London.
Single cortical neurons as deep artificial neural networks.
\emph{Neuron}, 109(17):2727--2739.e3, 2021.

\bibitem{sanchez2020}
A. Sanchez-Gonzalez, J. Godwin, T. Pfaff, R. Ying, J. Leskovec, and
P. Battaglia.
Learning to simulate complex physics with graph networks.
In \emph{Proceedings of the 37th International Conference on Machine
Learning}, 2020.

\bibitem{pfaff2021}
T. Pfaff, M. Fortunato, A. Sanchez-Gonzalez, and P. Battaglia.
Learning mesh-based simulation with graph networks.
In \emph{International Conference on Learning Representations}, 2021.

\bibitem{bengio2015}
S. Bengio, O. Vinyals, N. Jaitly, and N. Shazeer.
Scheduled sampling for sequence prediction with recurrent neural networks.
In \emph{Advances in Neural Information Processing Systems}, 2015.

\end{thebibliography}

\end{document}
